% BHU PYQ -- Manually transcribed & error-corrected
\documentclass[11pt,a4paper]{article}

\usepackage[a4paper,top=2.5cm,bottom=2.5cm,left=2.8cm,right=2.8cm,headheight=14pt]{geometry}
\usepackage{amsmath,amssymb}
\usepackage[T1]{fontenc}
\usepackage[utf8]{inputenc}
\usepackage{lmodern}
\usepackage{microtype}
\usepackage{fancyhdr}
\usepackage{enumitem}
\usepackage{hyperref}
\usepackage{lastpage}
\usepackage{parskip}

\hypersetup{colorlinks=true,urlcolor=black,linkcolor=black,
  pdftitle={CHB-602: Inorganic Chemistry-IV -- BHU PYQ},pdfauthor={Banaras Hindu University}}

\pagestyle{fancy}\fancyhf{}
\renewcommand{\headrulewidth}{0.4pt}\renewcommand{\footrulewidth}{0.4pt}
\fancyhead[L]{\small\textit{Banaras Hindu University}}
\fancyhead[R]{\small\textit{CHB-602: Inorganic Chemistry-IV}}
\fancyfoot[C]{\small Page \thepage\ of \pageref{LastPage}}

\newlist{parts}{enumerate}{2}
\setlist[parts,1]{label=(\alph*),leftmargin=*,itemsep=5pt,topsep=3pt}
\setlist[parts,2]{label=(\roman*),leftmargin=*,itemsep=3pt,topsep=2pt}
\newcommand{\pts}[1]{\hfill\textit{[#1]}}

\begin{document}

\begin{center}
  {\large\bfseries BANARAS HINDU UNIVERSITY}\\[2pt]
  {\normalsize B.Sc. (Hons.) Semester VI Examination, 2013-14}\\[6pt]
  \rule{0.8\linewidth}{0.5pt}\\[6pt]
  {\large\bfseries Paper No. CHB-602: Inorganic Chemistry-IV}\\[4pt]
  {\normalsize Time: 3 Hours\hfill Maximum Marks: 70}
\end{center}
\rule{\linewidth}{0.4pt}
\small
\noindent\textit{Instructions: Answer five questions in total, including Question~1 which is compulsory.}
\normalsize
\bigskip

\noindent\textbf{Question 1.} Answer all parts of the following: \pts{5 $\times$ 2 = 10}
\begin{parts}
  \item Why are low-spin complexes of $\text{Ni(II)}$ diamagnetic ($0.0\,\text{B.M.}$) whereas high-spin complexes are paramagnetic ($2.9\,\text{B.M.}$)?
  \item Draw the structure of the compound obtained by the reaction of uranyl nitrate with sodium acetate in acetic acid.
  \item Compare the electronic configurations of the first five actinides with those of the corresponding lanthanides.
  \item Describe the nature of the $\text{Co-C}$ bond present in methylcobalamin ($\text{Vitamin B}_{12}$).
  \item What is the role of the sodium-potassium pump in biological systems?
\end{parts}

\medskip\rule{\linewidth}{0.2pt}

\noindent\textbf{Question 2.}
\begin{parts}
  \item Why do $\text{Pu(III)}$ and $\text{Am(III)}$ show anomalous magnetic properties? Explain in terms of spin-orbit coupling. \pts{6}
  \item Explain the observed molar extinction coefficients ($\varepsilon$ in $\text{L}\,\text{mol}^{-1}\,\text{cm}^{-1}$) for the following complexes: \pts{8}
    \[ [\text{TiCl}_6]^{2-} \ (10000), \quad [\text{CoCl}_4]^{2-} \ (500), \quad [\text{Ti(H}_2\text{O)}_6]^{3+} \ (10), \quad [\text{MnBr}_4]^{2-} \ (4), \quad [\text{Mn(H}_2\text{O)}_6]^{2+} \ (0.02) \]
\end{parts}

\medskip\rule{\linewidth}{0.2pt}

\noindent\textbf{Question 3.}
\begin{parts}
  \item Describe the Gouy method for the determination of magnetic susceptibility of a complex. Compare it with the Faraday method. \pts{8}
  \item Account for the following: \pts{6}
  \begin{parts}
    \item $\text{KMnO}_4$ shows an intense violet color despite having a $\text{d}^0$ metal center.
    \item $\text{HgCl}_2$ is white while $\text{HgI}_2$ has a dark red color.
  \end{parts}
\end{parts}

\medskip\rule{\linewidth}{0.2pt}

\noindent\textbf{Question 4.}
\begin{parts}
  \item What are the selection rules for d-d electronic transitions? Describe the mechanisms of relaxation of these rules. \pts{6}
  \item Give a brief account of Charge Transfer (CT) transitions in the electronic spectra of coordination compounds. Differentiate between LMCT and MLCT with examples. \pts{8}
\end{parts}

\medskip\rule{\linewidth}{0.2pt}

\noindent\textbf{Question 5.}
\begin{parts}
  \item How is thorium extracted from monazite sand ore? Write down all reactions. \pts{6}
  \item How are transuranic elements separated from the spent fuel rods of a nuclear reactor (PUREX process)? \pts{8}
\end{parts}

\vfill
\rule{\linewidth}{0.4pt}
\begin{center}\small
  \href{https://bhu-exam.vercel.app/}{Click here to visit our website}\\[4pt]
  \href{https://me-aryan.vercel.app/}{Click here to visit our contributor}\\[4pt]
  \href{https://quashan-me.vercel.app/}{Click here to visit our contributor}
\end{center}

\end{document}
